\subsubsection*{定理 1.4.}
设 $X\subset\mathbb P^N$ 是域 $k$ 上的射影 Cohen--Macaulay 概形,
等维且维数为 $n\geq1$. 存在整数 $d_0$, 使对每个 $d\geq d_0$,
若 $i:X.H\longrightarrow X$ 表示一般的 $d$ 次超曲面截面的包含, 则
\begin{equation}\tag{1.4.0}
\mathfrak F_p\quad\text{作用于}\quad
\frac{H^{n-1}(X.H,\mathcal O_{X.H})}
     {i^*H^{n-1}(X,\mathcal O_X)}
\end{equation}
的作用非幂零.

\medskip
\noindent\emph{证明.}
如上, 可归结为 $k=\mathbb F_q$ 的情形, 并只需找到一个使 (1.4.0) 成立的 $H$:
其次数任取一个固定且充分大的值, 使 $\dim H^{n-1}(X.H,\mathcal O_{X.H})$
与 $H$ 无关; 见 (1.4.4). 扩大 $q$ 后, 可设 $X$ 的不可约分支 $X_i$ 几何不可约, 且
\begin{equation}\tag{1.4.1}
\operatorname{card}\bigl(X_i(\mathbb F_q)\bigr)\geq2,
\end{equation}
并且
\begin{equation}\tag{1.4.2}
\begin{aligned}
&\text{存在射影坐标 }X_0,\dots,X_N\text{ 使得}\\
&\qquad X\text{ 不交于线性子簇 }
X_0=\cdots=X_n=0.
\end{aligned}
\end{equation}
由 Serre 的消灭定理 ([1], 或 SGA~2 XII~1.4), 存在整数 $d_1$,
使对每个 $d\geq d_1$,
\begin{equation}\tag{1.4.3}
H^i(X,\mathcal O_X(-d))=0\qquad\text{若 }i\ne n=\dim X.
\end{equation}
因此, 若 $H$ 是次数 $d\geq d_1$ 的超曲面, 则 (1.4.3) 给出
\begin{equation}\tag{1.4.4}
\begin{aligned}
H^i(X,\mathcal O_X)&\xrightarrow{\ \sim\ }
 H^i(X.H,\mathcal O_{X.H}) &&\text{若 }i\leq n-2,\\[1ex]
H^{n-1}(X,\mathcal O_X)&\lhook\joinrel\longrightarrow
 H^{n-1}(X.H,\mathcal O_{X.H}) &&\text{为单射}.
\end{aligned}
\end{equation}

我们证明定理在 $d_0=\max(N+1,d_1)$ 时成立. 对每个 $d\geq d_0$,
需找到定义于 $\mathbb F_q$ 的 $d$ 次超曲面 $H$, 不含任何 $X_i$, 且使
\begin{equation}\tag{1.4.5}
\operatorname{tr}\!\left(\mathfrak F_q\ \middle|\
\frac{H^{n-1}(X.H,\mathcal O_{X.H})}
     {H^{n-1}(X,\mathcal O_X)}\right)\ne0.
\end{equation}
由迹公式 (XX~3.2.1), 模 $p$ 有
\begin{equation}\tag{1.4.6}
\begin{aligned}
\operatorname{card}\bigl((X.H)(\mathbb F_q)\bigr)
&\equiv\sum_{i=0}^{n-1}(-1)^i
 \operatorname{tr}\bigl(\mathfrak F_q\mid
 H^i(X.H,\mathcal O_{X.H})\bigr)\\
&\stackrel{(1.4.4)}{\equiv}
 \sum_{i=0}^{n-1}(-1)^i
 \operatorname{tr}\bigl(\mathfrak F_q\mid H^i(X,\mathcal O_X)\bigr)\\
&\qquad+(-1)^{n-1}\operatorname{tr}\!\left(\mathfrak F_q\ \middle|\
 \frac{H^{n-1}(X.H,\mathcal O_{X.H})}
      {H^{n-1}(X,\mathcal O_X)}\right).
\end{aligned}
\end{equation}
用 $a\in\mathbb Z/p\mathbb Z$ 表示下列量
\[
a=\sum_{i=0}^{n-1}(-1)^i
\operatorname{tr}\bigl(\mathfrak F_q\mid H^i(X,\mathcal O_X)\bigr).
\]
因此只需对每个 $d\geq d_0$ 找到定义于 $\mathbb F_q$、次数为 $d$ 且不含 $X$
的超曲面 $H$, 使得
\begin{equation}\tag{1.4.7}
\operatorname{card}\bigl((X.H)(\mathbb F_q)\bigr)\not\equiv a\pmod p.
\end{equation}

分两种情形. 若 $a\ne0$, 如 (1.3.8) 那样取 $\mathbb F_q$ 上使
$(X.H)(\mathbb F_q)$ 为空的 $H$ 即可. 这样的 $H$ 满足 (1.4.7) 且不含任何 $X_i$,
因为按假设 $\operatorname{card}(X_i(\mathbb F_q))\geq2$.

若 $a=0$, 只需找到使 $\operatorname{card}((X.H)(\mathbb F_q))=1$ 的 $H$.
此时 (1.4.7) 成立, 且由 $\operatorname{card}(X_i(\mathbb F_q))\geq2$ 得
$X_i\not\subset H$. 选取射影坐标 $X_0,\dots,X_N$, 使 $(1,0,\dots,0)$ 是 $X$ 的一个点.
取 $H$ 的方程为 $H(X_1,\dots,X_N)=0$, 其中 $H$ 是满足下式的齐次式 (参见 (1.2.2))
\begin{equation}\tag{1.4.8}
\alpha_1,\dots,\alpha_N\in\mathbb F_q,\quad
H(\alpha_1,\dots,\alpha_N)=0
\quad\Longrightarrow\quad
\alpha_1=\cdots=\alpha_N=0.
\end{equation}
于是 $H$ 在 $\mathbb P^N$ 中定义的超曲面在 $\mathbb F_q$ 上只有点 $(1,0,\dots,0)$;
更有 $X.H$ 在 $\mathbb F_q$ 上也只有此点. 定理得证.

\medskip
\noindent\emph{底本注.}
照录法文层恢复了扫描件中的 \emph{coordonnées}、(1.4.6) 中两个印出的等号,
以及 (1.4.8) 中的三点省略号; 它们均与继承的转录不同.
印本在 $d_o$ 和 $X_o$ 的零指标处使用字母形下标 $o$, 写作 \emph{une $H$ défini},
且计数点数的括号不规则. 校订法文层和英文层将其规范为 $d_0$、$X_0$、阳性一致形式
及通常的求值记号, 不改变论证.
